\section{Question, inherited objects, and the unresolved implication}
The preceding capacity theorem \cite{Capacity} completes a fixed-target divisor
problem. It proves that searching additional divisors at a fixed grade cannot
repair every exterior source. Its universal existence obligation is still open:
for a prescribed hard prime, produce an original exterior or middle state.
We do not assume that obligation in the present calculation.

The newly supplied ES--Zeta work \cite[\S\S1--3, equations (8)--(24)]{Input}
uses the \emph{literal} four roots $(p,x,y,z)$ and completes its signed inverse
by $\eta=\xi^{-1}$. Its monic equations give a finite flat rank-eight algebra.
The task here is to determine what happens to that particular algebra over
$\Z_p$, with its original coordinate lattice, rather than replacing it by an
arbitrary rank-eight vector space. The literal quartic is different from the
reciprocal cubic in \cite{Reciprocal}; the exact inversion and its derivative
twist are given in Section~\ref{sec:reciprocal}.

The Type-I/II identities are classical \cite[\S2, Propositions 2.2 and 2.6]{ET}.
Quadratic character restrictions are also antecedents; in particular compare
\cite[Corollary 1.3 and Appendix A]{BL}. All calculations needed for the local
and integral results below are proved explicitly. No historical-priority claim
is made for these specializations or for the standard local algebra tools.

\section{Original coordinates and the literal polynomial}
Fix a prime $p\equiv1\pmod{12}$. Thus $p\ge13$, and both $-1$ and $3$ are
squares modulo $p$. Set $K=\Q_p$, $\oo=\Z_p$, and normalize $\vp(p)=1$.
An original state retains
\[
 p/4<a<p/2,\quad R=4a-p,\quad u\mid a^2,
 \qquad E:R\mid4u+1,\qquad M:R\mid u+a.
 \tag{IC1}
\]
For $a=\prod_l l^{e_l}$ its original exponent vector is
$\beta_l=v_l(u)-e_l\in[-e_l,e_l]$. Let
\[
 g=\gcd(a,u),\qquad (h,r,s)=(g^2/u,u/g,a/g).
\]
Then $a=hrs$, $u=hr^2$, and $\gcd(r,s)=1$. The exact ordered returns are
\[
\begin{array}{ll}
 E:&\kappa=(pr+s)/R,\quad (x,y,z)=(a,hs\kappa,phr\kappa),\\
 M:&\lambda=(r+s)/R,\quad (x,y,z)=(a,phs\lambda,phr\lambda).
\end{array}                                                     \tag{IC2}
\]
Multiplication proves $4/p=1/x+1/y+1/z$. The ordered divisor inverses are
$a^2/(Ry-pa)$ and $pa^2/(Ry-pa)$, respectively. The middle complement
$u\mapsto a^2/u$ exchanges the last two denominators and is kept as an orientation,
not suppressed as a duplicate.

\begin{lemma}[Original character and unit facts]\label{lem:original}
Every factor in (IC2) is a $p$-unit. For each $d\mid a^2$,
$\Leg dR=\Leg dp$, with Jacobi symbol on the left. Consequently
\[
 E:\Leg hp=-1,\qquad M:\Leg{rs}p=-1.                      \tag{IC3}
\]
In the exterior case, $\kappa-r$ and $\kappa+r$ are $p$-units.
\end{lemma}
\begin{proof}
For an odd $l\mid a$, reciprocity and $R\equiv-p\pmod l$ give
$\Leg lR=\Leg{-1}l\Leg Rl=\Leg pl=\Leg lp$. If $2\mid a$, use
$R\equiv-p\pmod8$ and the supplementary law. Multiplication over the original
exponents proves the divisor identity. The E target is $-1/4$ modulo
$R\equiv3\pmod4$, so $\Leg up=-1$ and $u=hr^2$ gives the first assertion.
The M relation $r\equiv-s\pmod R$ gives the second.

The factors $h,r,s$ are below $p$. In E put
$\mathsf D=(4u+1)/R=(r+\kappa)/s$; the original identities give
$p\mathsf D+1=4hr\kappa$, so $\kappa$ is a unit too.
If $p\mid\kappa-r$, this identity gives $4hr^2\equiv1\pmod p$, contradicting
$\Leg hp=-1$. If $p\mid\kappa+r$, then $p\mid\mathsf D$ and
$4u+1=R\mathsf D$ gives the equally impossible nonresidue $u\equiv-1/4$.
For M, $r+s\le rs+1\le a+1<p$, so $\lambda<p$.
\end{proof}

The fixed, ordered root vector and scale are
\[
 (\ell_0,\ell_1,\ell_2,\ell_3)=(p,x,y,z),\qquad S=p+x+y+z,
\]
\[
 \mathcal H(T)=-S^{-1}\prod_{i=0}^3(T-\ell_i)
 =AT^4+T^3+BT^2+CT+D.                                   \tag{IC4}
\]
Here $D$ is a quartic coefficient, not the original exterior cofactor $\mathsf D$.
Vieta and the reciprocal identity give
\[
 A=-S^{-1},\quad C=5xyz/S,\quad D=-pxyz/S,
 \quad D=-pC/5,\quad B=-Ap^2-p-4C/(5p).                 \tag{IC5}
\]
These are the literal coefficient equations in \cite[\S2.1]{Input}; the prime
and every root label are retained. Define $d_i=\mathcal H'(\ell_i)$.

\begin{theorem}[Original prime-local root clusters]\label{thm:clusters}
The roots are distinct, $S\in\oo^\times$, and the following assertions hold.
For E write $(x,y,z)=(a,b,pc)$. Then $a,b,a-b,a+b$ are units,
$c\equiv1/4\pmod p$, and
\[
 (\vp(d_0),\vp(d_1),\vp(d_2),\vp(d_3))=(1,0,0,1),
 \qquad \vp\Disc(\mathcal H)=2.                         \tag{IC6}
\]
For M write $(x,y,z)=(a,pb,pc)$. Then $b,c,b-c,b-1,c-1$ are units,
\[
 4bc\equiv b+c\pmod p,\qquad\Leg{bc}p=-1,
\]
\[
 (\vp(d_0),\vp(d_1),\vp(d_2),\vp(d_3))=(2,0,2,2),
 \qquad \vp\Disc(\mathcal H)=6.                         \tag{IC7}
\]
In coefficient order $(A,B,C,D)$ the valuations are respectively
$(0,0,1,2)$ and $(0,1,2,3)$.
\end{theorem}
\begin{proof}
For E, $c=hr\kappa$ and $4c=p\mathsf D+1$. The remaining unit statements follow
from Lemma~\ref{lem:original} and $a\pm b=hs(r\pm\kappa)$.
Thus exactly one pair of roots, $p$ and $pc$, has difference of valuation one;
all other differences are units. In particular they are all distinct, and
$S\equiv a+b$ is a unit. The derivative and discriminant formulas prove (IC6).
The coefficient $B\equiv-ab/S$ is a unit, and (IC5) gives the other valuations.

For M, $b=hs\lambda,c=hr\lambda$. The unit and distinctness of $b,c$ follow from
$r\ne s$, $r,s<p$; equality would force $r=s=1$ and $R\mid2$.
The ES equation gives $4bc=pbc/a+b+c$, and (IC3) gives $\Leg{bc}p=-1$.
If $b=1\pmod p$, then $c=1/3$, making $bc$ a square since $\Leg3p=1$.
The same excludes $c=1$. The three roots $p,pb,pc$ therefore have pairwise
valuation-one differences; the singleton $a$ is separated by units. Now
$S\equiv a$, giving (IC7). In $B$ the leading term is
$-pa(1+b+c)/S$, and $1+b+c\equiv1+4bc$ cannot vanish, since $bc=-1/4$
would be a square. Hence $\vp(B)=1$. The remaining values follow from (IC5).
\end{proof}

In particular any two original E and M coefficient targets at the same prime
have distance exactly one in the maximum $p$-adic norm on $(A,B,C,D)$: the
$B$-coordinate is a unit on E and divisible by $p$ on M, while all coordinates
are integral. An infinitesimal target deformation preserving reduction is not
such a channel change. No comparison with a complex Gamma metric is inferred.

\section{The rank-eight generic algebra and its original integral order}
The supplied completion specializes, with its exact original basis, to
\[
 \OO=\oo[T]/(\mathcal H),\qquad
 \SSS=\OO[\eta]/(\eta^2-\mathcal H'(T)).                 \tag{IC8}
\]
The leading coefficient $A$ is a unit, so both equations can be made monic.
Thus $\SSS$ is free over $\oo$ on
\[
 1,T,T^2,T^3,\eta,\eta T,\eta T^2,\eta T^3.             \tag{IC9}
\]
After inverting $\eta$ it is the original signed inverse, with
$\xi=\eta^{-1}$ and $\xi^2\mathcal H'(T)=1$.
The full physical Fabel coordinates, inherited from \cite[equation (21)]{Input}, are
\[
 \left(\eta^{-1},-T-i\eta,
 A\eta^3+2T\eta+3i\eta^2,
 7iT^2\eta+(B-17T+AT^2)\eta^2-13i\eta^3-2A\eta^4\right).       \tag{IC10}
\]
Either embedding of $i$ in $K$ may be chosen and is kept fixed. No complex
period coefficient is presumed to have a $p$-adic descent.

Over $K$, root evaluation gives exactly
\[
 \SSS_K\simeq\prod_{i=0}^3K[\eta_i]/(\eta_i^2-d_i).     \tag{IC11}
\]
There is no infinity factor: $A\ne0$ and all four roots are finite.
This is not the seven-dimensional reciprocal-cubic algebra.

\begin{theorem}[Generic local signed types]\label{thm:fields}
An E state has either
\[
 K^4\times K_{\rm ram}^2
 \quad\text{or}\quad K_{\rm unr}^2\times K_{\rm ram}^2,
                                                                  \tag{IC12}
\]
and consequently four or zero $K$-rational signed inverse points.
An M state has either
\[
 K^8\quad\text{or}\quad K^4\times K_{\rm unr}^2,          \tag{IC13}
\]
and consequently eight or four such points. Repeated field factors retain their
root labels; they are not identified.
\end{theorem}
\begin{proof}
For E the two valuation-one derivatives have square ratio: after dividing by
$p$ their ratio is $-1\pmod p$. They define the same ramified quadratic extension.
At the other two roots the derivative residues are
\[
 -\frac{a^2(a-b)}{a+b},\qquad -\frac{b^2(b-a)}{a+b}.
\]
Their ratio is a square, so either both split or both define the unramified
quadratic extension. This proves (IC12).

For M, divide the three cluster derivatives by $p^2$. Their unit residues are
\[
 (1-b)(1-c),\quad(b-1)(b-c),\quad(c-1)(c-b).              \tag{IC14}
\]
Their product is the negative of a nonzero square. Since $-1$ is a square,
either all three or exactly one are squares. At the singleton,
$d_1\equiv-a^2$ is a square. This proves (IC13).

For the local facts, a square unit modulo the odd prime has unique compatible
lifts after a sign is chosen: if $z_n^2\equiv d\pmod{p^n}$, set
$z_{n+1}=z_n+p^nk$, with
$k\equiv-(z_n^2-d)/p^n\,(2z_n)^{-1}\pmod p$.
A nonsquare unit has irreducible separable quadratic reduction and gives the
unramified quadratic extension. An odd-valuation class gives a ramified
quadratic. These are also the standard constructions in
\cite[Chapter 7, Newton's lemma and local extensions]{Milne}.
\end{proof}

At the single hard prime $p=1009$, all four alternatives occur. The respective
original $(a,u)$ and ordered denominators are
\[
\begin{array}{c|c|c}
\text{channel, number of }K\text{-points}&(a,u)&(x,y,z)\\\hline
E,0&(253,11)&(253,87032,3818056)\\
E,4&(253,23)&(253,86020,7890380)\\
M,8&(253,11)&(253,2042216,88792)\\
M,4&(253,23)&(253,1021108,92828).
\end{array}                                                   \tag{IC15}
\]
These are four independently checked original states. In particular a
$K$-rational point on the literal signed fibre is not a necessary condition for
an ES witness. Every row still has eight geometric characteristic-zero points.

\section{Exact normalization, conductors and invariant factors}
Put $k_i=\lfloor\vp(d_i)/2\rfloor$ and
\[
 \theta_i=\eta_i/p^{k_i},\qquad
 \delta_i=d_i/p^{2k_i},\qquad
 N_i=\oo[\theta_i]/(\theta_i^2-\delta_i),\qquad
 \NN=\prod_iN_i.                                         \tag{IC16}
\]
Each $\delta_i$ is a unit or has valuation one. Thus $N_i$ is the maximal
integral order in its generic factor. For a field factor, integrality of
$a+b\theta_i$ first gives $a\in\oo$ by its trace. Its norm then gives
$2\vp(b)+\vp(\delta_i)\ge0$, forcing $b\in\oo$ because $\vp(b)$ is integral.
For a split unit factor use the two values $a\pm b\sqrt{\delta_i}$ and the
unit $2\sqrt{\delta_i}$. This also proves maximality without an unspecified
integral basis.

The precise normalization map is
\[
 f(T)+\eta g(T)\longmapsto
 \bigl(f(\ell_i)+p^{k_i}g(\ell_i)\theta_i\bigr)_{i=0}^3,
 \qquad \deg f,\deg g<4.                                \tag{IC17}
\]
It is injective: over $K$, evaluate at all distinct labelled roots.
All its inverse coefficients are provided by the corresponding Lagrange
polynomials, so its integral image can be tested without changing the lattice.

\begin{theorem}[Complete original integral defect]\label{thm:index}
For E,
\[
 \NN/\SSS\simeq(\oo/p)^2,
 \qquad[\NN:\SSS]=p^2,
 \qquad\mathfrak f_E=pN_0\times N_1\times N_2\times pN_3.       \tag{IC18}
\]
For M,
\[
 \NN/\SSS\simeq(\oo/p)^2\oplus(\oo/p^2)^2\oplus\oo/p^3,
 \qquad[\NN:\SSS]=p^9,
\]
\[
 \mathfrak f_M=p^3N_0\times N_1\times p^3N_2\times p^3N_3.      \tag{IC19}
\]
Here $\mathfrak f=\{n\in\NN:n\NN\subseteq\SSS\}$ is the conductor ideal in
this finite algebra. It is not the analytic weighted conductor of the Zeta
programme. In the eight-dimensional original basis, the Smith exponents are
\[
 E:(0,0,0,0,0,0,1,1),\qquad
 M:(0,0,0,1,1,2,2,3).                                   \tag{IC20}
\]
\end{theorem}
\begin{proof}
Unit differences split off every noncolliding root integrally, by the polynomial
CRT. For E, the remaining root lattice on the pair $p,pc$ is exactly
\[
 L_2=\{(z_0,z_3)\in\oo^2:z_0-z_3\in p\oo\}.
\]
A linear interpolation proves sufficiency, since $p-pc$ is $p$ times a unit.
Both coefficients $f$ and $g$ in (IC17) give a copy of $L_2$: here $k_i=0$.
This proves the quotient and exponents. A normal element supported at just one
of these two roots is in the original order precisely when its coefficients
of $1$ and $\theta_i$ both lie in $p\oo$. The resulting $pN_i$ is an ideal and
is maximal with this property, giving (IC18).

For M, order the cluster as $p,pb,pc$ and use the Newton polynomial basis
$1,T-p,(T-p)(T-pb)$. Evaluation gives diagonal factors $1,p,p^2$ times units,
and integral row operations remove the lower entries. Thus its root lattice
$L_3$ has quotient $\oo/p\oplus\oo/p^2$. The two coefficients in (IC17)
give $L_3\oplus pL_3$. Their exponents are $(0,1,2)$ and $(1,2,3)$; the
singleton contributes two zeros. This proves (IC19)--(IC20).

For the conductor, a vector supported at one of the three roots has a unique
Lagrange interpolation of degree at most two. Its leading coefficient has
denominator of exact order $p^2$. Hence a normal value $a+b\theta_i$ supported
there belongs to $\SSS$ precisely when $a\in p^2\oo$ and $b\in p^3\oo$.
For it to lie in the conductor, multiplication by the unit $\theta_i$ must also
belong to the order. This forces $a,b\in p^3\oo$, and that condition is
sufficient because $1,\theta_i$ is the full normal basis. The singleton is
already maximal. This proves the displayed conductor.
\end{proof}

The root orders alone have indices $p$ (E) and $p^3$ (M), with Smith exponents
$(0,0,0,1)$ and $(0,0,1,2)$. A further exact check is the trace discriminant:
\[
 \vp\Disc(\SSS)=2\vp\Disc(\OO)+\sum_i\vp(4d_i)
 =6\ (E),\quad18\ (M).
\]
The normal algebra has discriminant valuation two (E) or zero (M). The square
of the change-of-basis determinant again gives index exponents two and nine.
Neither determinant is a density or positivity assertion.

\begin{corollary}[The original actions on the finite defect module]
Let $Q=N/\widehat S$, a module rather than a quotient ring. Multiplication by the
original $T$ and $\eta$ descends to it. Their exact nilpotence indices are
\[
 E:(1,2),\qquad M:(3,3),                                \tag{IC20a}
\]
respectively. Thus the finite quotient retains nonzero action, even though
rational extension removes it.
\end{corollary}
\begin{proof}
The quotient is supported only at the colliding roots. For E, $T$ is divisible
by $p$ there and $\eta^2=d_i$ is also divisible by $p$, so the conductor kills
both actions. Multiplying a supported normal constant by $\eta$ still has a
coefficient not divisible by $p$ and is nonzero in $Q$.
For M, both $T^3$ and $\eta^3$ have a factor $p^3$ at each colliding normal
factor. But their squares, applied to a supported normal $\theta_i$, have a
$\theta_i$-coefficient of exact valuation two. The membership criterion in the
proof of Theorem~\ref{thm:index} requires three, so both squared actions remain
nonzero. All the normal root quotients are units, precluding cancellation.
\end{proof}

\paragraph{Complete inverse coordinates.}
For a normal vector $a_i+b_i\theta_i$, interpolate the four values $a_i$ to
obtain $f(T)$, and the values $b_i/p^{k_i}$ to obtain $g(T)$. Membership in the
original order means all eight coefficients of $f,g$ are in $\oo$. Thus (IC17)
has an explicit partial integral inverse and a total $K$-linear inverse.
The quotient in (IC18) or (IC19) is retained when the integral inverse fails.
There is no $\oo$-linear splitting of its surjection from $\NN$, since that
would map a nonzero torsion module into a torsion-free one.

\section{Boundary fibres and a sharp finite lifting obstruction}
At E, the colliding factor of the special fibre is
\[
 \F_p[\epsilon,\eta]/(\epsilon^2,\eta^2-2g_0\epsilon)
 \simeq\F_p[\eta]/(\eta^4),\qquad g_0=-ab/S\ne0.         \tag{IC21}
\]
At M it is
\[
 \boxed{\F_p[\epsilon,\eta]/(\epsilon^3,\eta^2-3\epsilon^2).} \tag{IC22}
\]
These follow from the actual reduced quartics, respectively
$-S^{-1}T^2(T-a)(T-b)$ and $T^3(1-T/a)$, using the supplied
multiple-root algebra \cite[equation (23)]{Input}. Their dimensions are four
and six. Multiplication by $\epsilon$ has blocks $(2,2)$ or $(3,3)$.
Multiplication by $\eta$ has blocks $(4)$ or $(4,2)$: on
$1,\epsilon,\epsilon^2,\eta,\epsilon\eta,\epsilon^2\eta$ in (IC22), its
successive ranks are $6,4,2,1,0$. Thus (IC22) is not a single length-six chain.

The normalization map modulo $p$ has kernel dimension two on the E collision
block and five on the M collision block. For E, $T$ becomes zero in each normal
factor modulo $p$, while $\eta$ gives its two labelled nonzero nilpotents. For M,
both $T=p\zeta$ and $\eta=p\theta$ become zero in the normal special fibre.
These are the kernels dictated by (IC20), not a deletion of the original labels.

\begin{theorem}[The added E point stops at the third digit]\label{thm:lift}
The point $(T,\eta)=(0,0)$ of (IC21) has exactly $p$ lifts solving both original
equations modulo $p^2$, and has no lift modulo $p^3$.
\end{theorem}
\begin{proof}
Write $\mathcal H(T)=U(T)(T-p)(T-pc)$, with
$U(T)=-(T-a)(T-b)/S$ a unit near zero. A putative lift modulo $p^2$ is
$T=pw,\eta=pv$, with $w,v\in\F_p$. The first equation is automatic. The
second gives
\[
 0\equiv pU(0)(2w-1-c)\pmod{p^2},
\]
so $w=(1+c)/2=5/8\pmod p$. The $p$ values of $v$ are all possible.
For any higher lift of that $T$,
\[
 \frac{\mathcal H(T)}{p^2}\equiv
 -U(0)\frac{(1-c)^2}{4}
 =\frac{9ab}{64S}\not\equiv0\pmod p.                   \tag{IC23}
\]
Thus the root equation already fails modulo $p^3$, independently of $\eta$.
\end{proof}

The finite flat completion is proper, but this boundary point has no integral
$\oo$-section. The ramified generic points remain over their actual quadratic
extensions. At $p=1009$ in the first row of (IC15), the entire generic signed
fibre has no $K$-point despite an original ES solution and a rational special
boundary point. No smooth Hensel hypothesis is available at $\eta=0$.

For M, the collision block instead has the exact integral normalization chart
\[
 G(Z)=(Z-1)(Z-b)(Z-c),\qquad
 G(Z)=0,\quad \theta^2=\frac{a-pZ}{S}G'(Z),              \tag{IC24}
\]
with $T=pZ,\eta=p\theta$. This is finite etale of rank six, since the reduced
roots are distinct and the right side is a unit there. The inverse divides
exactly by those powers of $p$. The character product in (IC14) proves that
at least one cluster root has two $K$-rational signed lifts. Thus the M boundary
point does lift to $\oo$, although its full multiplicity and the index $p^9$
are still retained.

The original physical inverse (IC10) has pole $\vp(\xi)=-1/2$ on the E cluster
and $\vp(\xi)=-1$ on the M cluster. Inverting $\eta$ in (IC21) or (IC22)
annihilates that local algebra; it does not turn the added state into a finite
original affine inverse. The supplied complex metric limit remains the one in
\cite[equations (21)--(22)]{Input}; no identification of that complex norm with
this nonarchimedean valuation is asserted.

\section{The generic prime-sign symmetry has an integral defect}
The distinguished-root projector is
\[
 e_p(T)=\frac{\prod_{i=1}^3(T-\ell_i)}{\prod_{i=1}^3(p-\ell_i)}.
                                                               \tag{IC25}
\]
Its values are $(1,0,0,0)$. Its coefficients require an exact denominator $p$
in E and $p^2$ in M, since its numerator is monic and the denominator has
those respective valuations. This is also the supplied $g_p/d_p$ projector
\cite[equations (9)--(11)]{Input}.

Among the supplied four generic deck sign maps, the prime-only flip is
\[
 T\longmapsto T,\qquad \eta\longmapsto(1-2e_p(T))\eta.
                                                               \tag{IC26}
\]
Uniqueness in the original free basis (IC9) proves that it does not preserve
$\SSS$. Nor does its negative, the denominator-only flip. The identity and the
simultaneous sign $\eta\mapsto-\eta$ do preserve it. All four extend to the
normalization (IC16), where their action is an independent sign on the labelled
quadratic factors. This conclusion concerns those four source deck maps, not
all automorphisms of a specialized split algebra.

\section{Literal channel interpolation and reciprocal inversion}
Suppose actual E and M states are both supplied at the same $p$. Denote their
ordered literal roots by $x_i$ and $y_i$. The exact root-algebra maps are
\[
 F(T)=\sum_i y_i\frac{\prod_{j\ne i}(T-x_j)}{\prod_{j\ne i}(x_i-x_j)},
 \qquad
 G(T)=\sum_i x_i\frac{\prod_{j\ne i}(T-y_j)}{\prod_{j\ne i}(y_i-y_j)}.
                                                               \tag{IC27}
\]
Evaluation proves both inverse compositions on the respective quotient rings.
Under their common label coordinates in $K^4$, the integral orders obey
\[
 \OO_M\subset\OO_E,
 \qquad[\OO_E:\OO_M]=p^2.                              \tag{IC28}
\]
Indeed the three M cluster values are congruent modulo $p$, in particular at
the E pair of labels $0,3$. The already proved indices give the ratio.
Thus $F$ is integral. In the reverse interpolation, the summand at label $2$
has unit output $x_2$ and a denominator of valuation two; the other cluster
outputs are divisible by $p$. Consequently
\[
 \boxed{\min_j\vp(G_j)=-2,\qquad
 p^2G(T)\bmod p=c_0T^2(T-a_M),\quad c_0\ne0.}           \tag{IC29}
\]
This proves the exact lost factor, not just an upper denominator estimate.
At $p=1009$, the pair of $(a,u)=(253,11)$ states in (IC15) gives
$c_0=306$ and $p^2G\equiv275T^2+306T^3\pmod{1009}$.

The full signed generic algebras need not be isomorphic over $K$: E has ramified
field factors and M has none. The surviving comparison is explicit after a
specified extension. Adjoin square roots $\delta_i^2=d_i^M/d_i^E$ and send
$\eta_M$ to $\delta(T)\eta_E$, where $\delta$ is the labelled interpolation
of the $\delta_i$. The inverse uses their reciprocals. The extension has degree
at most four: for odd $p$, unit classes are square or a fixed nonsquare unit,
and the only additional parity is that of $p$. The original integral images
must still be tested by (IC17); they are not made unit lattices by this field
isomorphism.

This is applicable to actual source-changing returns, for example the preceding
capacity certificate \cite[sharp finite repair]{Capacity} at $p=825241$:
\[
 E:(a,R,u,h,r,s,\kappa)=(206321,43,2240439739,19,10859,1,208402140),
\]
\[
 M:(a,R,u,h,r,s,\lambda)=(217170,43439,25,1,5,43434,1).
\]
The latter is supplied there by the noncanonical original cofactor word $25$,
not by interpolation in (IC27). The current certificate transports both actual
states through (IC18)--(IC29), preserving their different factor inventories.
No M state is assumed to exist in order to use (IC27) as an existence proof.

\subsection{The reciprocal model changes the lattice and the signed square class}
\label{sec:reciprocal}
Keep all four roots while applying $z_i=p/\ell_i$. On the literal root algebra,
\[
 \boxed{I(T)=-\frac pD(AT^3+T^2+BT+C)=p/T.}              \tag{IC30}
\]
Its exact coefficient denominator is $p$ in E and $p^2$ in M. The transformed
monic polynomial is $Z^4L(p/Z)/L(0)$, where $L=\prod(T-\ell_i)$.
If $N=\prod_i\ell_i$ and the normalized reciprocal quartic is
$\widetilde{\mathcal H}=-\frac15\prod(Z-z_i)$, differentiation gives
\[
 \boxed{\widetilde{\mathcal H}'(z_i)=
 \omega\frac{\mathcal H'(\ell_i)}{\ell_i^2},\qquad
 \omega=-\frac{p^3S}{5N}.}                              \tag{IC31}
\]
The four reciprocal roots sum to five; the coefficient $-1/5$ is therefore
fixed, not fitted. The signed comparison needs the recorded quadratic twist
$\sqrt\omega$. Its valuation is one in E and zero in M.

The reciprocal discriminant valuation is the literal value plus
$12-6\vp(N)$, hence two in E and zero in M. The reciprocal root order has
index $p$ in E and is maximal in M. In E its colliding labels are now $1,2$,
not $0,3$. Removing the separately marked root $z_0=1$ yields the preceding
reciprocal cubic, whose signed derivatives acquire the further factors
$-(z_i-1)/5$ from the augmented quartic. This states the maps and square classes
rather than conflating the earlier seven-point and current eight-point fibres.

\section{A single missing prime passes any fixed finite integrality window}
The supplied real separation estimates \cite[\S1.2, equations (2)--(4)]{Input}
are necessary at an integral witness. The next result gives a precise scope
test for treating them, together with a finite local completion, as sufficient.
It retains both numerical channel gates and the entire allowed numerical range
of the proposed integer $u$. The one failed condition is original square-divisor
availability, at one explicitly designated prime.

\begin{theorem}[A joint, single-prime source obstruction]\label{thm:controls}
For every fixed finite set of rational primes $\Sigma$, there are infinitely
many hard primes $p$ and integers $a,R,u$ such that
\[
 p/4<a<p/2,\quad R=4a-p=3,\quad0<u<a^2,
 \quad R\mid4u+1,\quad R\mid u+a,                       \tag{IC32}
\]
yet $u\nmid a^2$. There is just one extra prime occurrence in $u$, at a
specified prime $l\notin\Sigma\cup\{p\}$. In particular
\[
 a^2/u\in\Z_q\quad(q\in\Sigma\cup\{p\}),
 \qquad v_l(a^2/u)=-1.
\]
The literal E and M rational denominator returns from these same numerical
coordinates satisfy the ES identity and all the following conditions:
\begin{enumerate}
\item their only nonintegral denominator has reduced denominator exactly $l$;
\item at $p$ they have the respective E and M local signatures, normalization
indices, conductors and invariant factors proved above;
\item their four positive literal roots have successive gaps greater than one;
\item with $s_p=(p^2+3p)/4$ and $B_p=p+s_p(s_p+1)$, both satisfy
\[
 S<B_p,\quad\max(x,y,z)<s_p(s_p+1)/3,\quad
 \Disc(\mathcal H)\ge144/B_p^6,\quad
 2/S\le|\mathcal H'(\ell_i)|\le S^2;
\]
\item the same prime $p$ has separately supplied original integer E \emph{and}
M states. The failed source in (IC32) is not an ES counterexample.
\end{enumerate}
\end{theorem}
\begin{proof}
Choose a prime $l\equiv11\pmod{12}$ outside $\Sigma$. Dirichlet's theorem
\cite[Theorem 18.1]{Dirichlet} permits this choice, and then supplies infinitely
many primes in the reduced CRT progression
\[
 p\equiv1\pmod{840},\qquad p\equiv-1\pmod l.
\]
The moduli are coprime. Take $p>4l$ and set
\[
 a=(p+3)/4,\quad R=3,\quad u=al,\quad
 h=a/l,\quad r=l,\quad s=1,\quad
 \kappa=(pl+1)/3,\quad\lambda=(l+1)/3.                  \tag{IC33}
\]
Both quotients are positive integers, and $l<a$ gives $0<u<a^2$.
Modulo three, $p\equiv a\equiv1$ and $l\equiv2$, so both gates hold.
But $a\equiv1/2\pmod l$ and $\gcd(a,u)=a$. Thus $h=a/l$ is nonintegral with
exact denominator $l$. Every prime of $a$ occurs in $u$ with its original
exponent; the only extra coordinate is $v_l(u)=1$, while $v_l(a)=0$.
In the centered original box this is exactly $\beta_l=1$ outside $[0,0]$.
All the equations $a=hrs$, $u=hr^2$, and $\gcd(r,s)=1$ hold, but the positive
rational grade $h$ is not an allowed original integer grade.

The two explicit returns are
\[
 E_{\rm rat}:\quad(a,a\kappa/l,pa\kappa),\qquad
 M_{\rm rat}:\quad(a,pa\lambda/l,pa\lambda).              \tag{IC34}
\]
Multiplication proves their reciprocal identities. Since $p,a,\kappa,\lambda$
are all units at $l$, the middle displayed denominator has exact reduced
denominator $l$ and the other two are integers. Every declared finite place
therefore passes the denominator test and the original divisor valuation test.

At the original prime, $p\nmid a,l,\kappa,\lambda$ and
$(l/p)=(p/l)=(-1/l)=-1$. Also $a\equiv3/4\pmod p$ is a square. Thus the same
unit and character calculations as Theorem~\ref{thm:clusters} apply to (IC34):
in the first return $4a\kappa\equiv1$, while the second has normalized root
product $a^2\lambda^2/l$ a nonsquare. The reduction patterns, local field
alternatives and all integral normalization formulas follow exactly. This
verifies local algebra for rational targets; it does not insert a missing
integer factor into the original box.

For both triples write their largest denominator as $z=s_p q$,
where $q=\kappa$ or $\lambda$ and $s_p=pa$. Since $l<a$,
$\kappa=(pl+1)/3<(s_p+1)/3$, and the same bound holds for $\lambda$.
Hence $z<s_p(s_p+1)/3$. Their other denominator is $z/(pl)$ in E and $z/l$ in M.
It is greater than $pa/3>p+1$; also $p-a>1$ and $z-y>1$. Thus the ordered root
list $a,p,y,z$ has successive gaps greater than one, and
$S=p+a+y+z<p+3z<B_p$. Four such roots have Vandermonde modulus at least twelve.
The product of three differences at a selected root is at least two and at most
$S^3$. The claimed discriminant and derivative bounds follow from
$\mathcal H=-S^{-1}\prod(T-\ell_i)$.

Finally retain the genuine endpoint data
\[
 R_E=l,\quad a_E=u_E=(p+l)/4,\quad
 (h_E,r_E,s_E,\kappa_E)=(a_E,1,1,(p+1)/l).
\]
They are integral and return $(a_E,a_E\kappa_E,pa_E\kappa_E)$.
There is also a genuine middle state: put
$j=(l+1)/4$, $c=(p+1)/l$, then
\[
 (h_M,r_M,s_M,\lambda_M)=(j,1,c,1),\quad a_M=jc,\quad
 R_M=c+1,\quad u_M=j.
\]
Its identity is $p=lc-1=4jc-c-1$. Since $l\ge11$ and $c>4$, both first-half
inequalities hold. This gives the last assertion without assuming ES.
\end{proof}

For $\Sigma=\{2,3,5,7,11\}$, choose $l=23$ and the hard prime $p=35281$.
The numerical source is
\[
 a=8821,\quad R=3,\quad u=202883,
 \quad(h,r,s,\lambda,\kappa)=(8821/23,23,1,8,270488).
\]
The two nonintegral denominator triples are
\[
 \boxed{E_{\rm rat}=(8821,2385974648/23,84179571556088),}
\]
\[
 \boxed{M_{\rm rat}=(8821,2489709608/23,2489709608).}      \tag{IC35}
\]
The E target has zero rational local signed points and the M target has four.
Their local normalization and ramification types are different. The only forbidden original factor is
$23\mid u$, while $23\nmid a=8821$. A second complete example at $p=54601$ is
included in the certificate. Trial division proves these displayed primes.

This theorem does not concern all places simultaneously: the omitted place is
explicit. Nor does it exclude an adaptive test of every possible denominator
prime. It proves that a fixed finite collection of local observations, even
with the original numerical ranges, both channel congruences and the stated
real bounds, is not original square-divisor membership. For the broader
arithmetic-approximation context, compare \cite[Theorems 1.1, 1.8--1.9]{BL}.
The separate companion fragment \texttt{crt\_control\_note.tex} preserves an
additional exact CRT construction with three larger-prime control targets;
those targets and their full equations are also independently checked.

\section{Verification scope, information loss, and continuation}
The full replay covers every original first-half box for the 99 primes
$p\equiv1\pmod{12}$ through 3000: 868359 square divisors and 4505 labelled
states, of which 2451 are E and 2054 are oriented M. It compares all derivative
valuations, characters, coefficient identities and local point alternatives.
The deep lattice calculations expand every original normalization matrix at
the four $p=1009$ examples. The separate checker uses all minors for the
invariant factors and independent Lagrange inversion, rather than the first
program's elimination. Both test the direct $p^2/p^3$ boundary equations, the
actual capacity crossing, and every supplied rational control.

The main verifier is standard-library Python. The separate implementation imports
none of it or of the preceding work. Its original source enumeration is by
primitive coprime $r,s$ and $h=a/(rs)$, rather than by divisors of $a^2$.
The prime domains, original vectors, ordered states and finite computation
ranges are all recorded in JSON. Checks of a finite prefix are not a proof
that the original source is nonempty at every prime.

No source coefficient is created by normalization. The exact sequence
$0\to\SSS\to\NN\to\NN/\SSS\to0$ retains all the finite torsion coordinates
in (IC18)--(IC20). Taking a quotient, a rational extension, or a reduction modulo
$p$ is accompanied by its calculated kernel. A receiving zero is not relabelled
as absence of an original divisor word, and an added nilpotent boundary point
is not promoted to an original E/M word. The countermodels in
Theorem~\ref{thm:controls} make that latter distinction quantitative.

The remaining Turn~7 task is still to force a positive original integer
coefficient for each prescribed hard prime. The completed same-grade capacity
classification remains available, and its actual source-changing maps are
compatible with the lattice comparison here. Finite rank, properness, local
normalization and the displayed approximation bounds do not supply the initial
arithmetic state by themselves. No Turn~8 witness guarantee, universal ES
proof, Lean build or independent mathematical review is claimed.
